\section*{Tóm tắt}
\phantomsection
\addcontentsline{toc}{section}{Tóm tắt}

Báo cáo trình bày có hệ thống bài toán phân loại đa lớp, từ phát biểu hình thức, các hàm mất mát và giới hạn tổng quát hóa đến những nhóm thuật toán tiêu biểu. Nội dung tập trung vào hai hướng chính: các mô hình học trực tiếp trên không gian nhãn đa lớp và các phương pháp quy đổi thành nhiều bài toán nhị phân như One-vs-All, One-vs-One và Error-Correcting Output Codes. Bên cạnh nền tảng lý thuyết, báo cáo phân tích các tiến bộ gần đây trong phân loại đa nhãn cực lớn, bao gồm mô hình dựa trên cây nhãn, biểu diễn nhãn, Transformer và chiến lược tạo danh sách ứng viên. Cuối cùng, báo cáo đề xuất một thiết kế demo thực nghiệm nhằm so sánh các phương pháp theo độ chính xác, macro-F1, chi phí huấn luyện và khả năng mở rộng.

\noindent\textbf{Từ khóa:} phân loại đa lớp, phân loại đa nhãn cực lớn, One-vs-All, One-vs-One, ECOC, cây nhãn, Transformer.

\subsection*{Đóng góp chính}
\phantomsection
\addcontentsline{toc}{subsection}{Đóng góp chính}

\begin{itemize}[leftmargin=2em]
  \item Hệ thống hóa cơ sở lý thuyết của phân loại đa lớp, bao gồm ký hiệu, hàm mất mát, đánh giá và các giới hạn tổng quát hóa.
  \item So sánh các nhóm thuật toán trực tiếp và kết hợp theo số mô hình, chi phí tối ưu, cơ chế dự đoán và hạn chế thực nghiệm.
  \item Phân tích các hướng nghiên cứu hiện đại cho bài toán nhãn cực lớn, đặc biệt là cây nhãn, nhúng nhãn, mô hình Transformer và cơ chế shortlist--rerank.
  \item Đề xuất quy trình demo thực nghiệm để đánh giá đồng thời chất lượng dự đoán, tính ổn định và khả năng mở rộng của các phương pháp.
\end{itemize}

\clearpage